\Needspace{8\baselineskip}
\section*{Question 2}
\noindent\textit{English main text: \QtwoWords\ words.}
\subsection*{(1) Involved modules}
The proposed delivery system uses all seven parts of the Q1 framework: hardware and sensors, perception and data processing, positioning and SLAM, mission planning, motion planning, modelling and control, and GCS/HMI with communications (\citealp[p. 21]{Chen2022}). The unmanned ground vehicle (UGV) provides a mobile ground station and recovery deck. I assume bounded vehicle speed and acceleration, a clear, approximately level deck, adequate payload capacity and battery reserve, and a balcony tray with clear overhead access and a registered visual identifier. These conditions define a practical setting for integrating the methods below.

\subsection*{(2) Required features and example methods}
Before departure, the GCS supplies the delivery manifest and operating constraints. A behaviour tree organises execution through conditions, actions and fallback branches (\citealp{Colledanchise2018}). In this proposal, readiness checks cover propulsion, sensors, communication, payload retention and return energy. The tree advances when observations confirm completion, allowing failed searches or infeasible trajectories to trigger bounded retries or return. This connects mission decisions to feedback from the other modules throughout the flight.

Reliable state estimation is essential from the moving launch onward. An EKF-based visual-inertial estimator combines IMU predictions with camera corrections (\citealp{Mourikis2007}), while reliable GNSS updates and local SLAM mapping are added in this integration (\citealp{Cadena2016}). Calibrated observations of known deck-marker points provide relative pose through PnP (\citealp{Lepetit2009}). Timestamped estimates are expressed in a common frame, allowing the UAV to maintain the deck's horizontal velocity while climbing clear and then transition into outbound flight.

During transit, I combine A* graph search (\citealp{Hart1968}) with constrained minimum-snap polynomial planning (\citealp{Mellinger2011}; \citealp{Richter2016}). The resulting timed references must satisfy obstacle clearance, speed and thrust limits. A rigid-body and rotor-thrust model supports cascaded P/PID control: position and velocity loops request thrust and attitude, while attitude and rate loops request moments. Velocity and acceleration feedforward assist tracking, and allocation produces four feasible rotor commands (\citealp{PX4}). New obstacle observations require the planner to revise the trajectory.

Near the apartment building, a planned facade scan maintains visibility and clearance while camera and range measurements locate candidate balconies, people and obstacles. YOLO provides the detection method (\citealp{Redmon2016}), with training adapted to balcony imagery. The system then checks the building, floor and apartment specified in the manifest against the registered identifier. This distinction matters because recognising a balcony does not establish the intended recipient. Uncertain identity or an obstructed approach prevents delivery.

Once identity and clearance are confirmed, the UAV hovers above the accessible tray. I propose a short winch that lowers the parcel while load or contact sensing checks that the tray supports it. Release is permitted only after this condition is met, and camera evidence confirms deposition before the mechanism retracts. These physical checks establish delivery beyond merely issuing a release command. The controller then updates its mass-dependent hover thrust to account for the removed payload.

The outcome must also reach the truck. MAVLink carries timestamped vehicle-motion estimates, together with battery and system-status telemetry (\citealp{MAVLink}). The proposed application sends a delivery record containing the target identity, outcome, time and evidence, with acknowledgement and retries to verify reception. Report reception and physical delivery remain separate completion conditions; a protocol command acknowledgement alone establishes neither. High-rate stabilisation stays onboard, so flight control does not depend on each telemetry exchange.

Recovery closes the mission by updating the moving rendezvous and tracking deck-relative position and velocity (\citealp{Borowczyk2017}). My design permits descent only with alignment, low horizontal relative speed and reliable tracking. During descent, it continues to match the deck's horizontal motion, reducing vertical relative speed towards zero as it reaches touchdown. Stale updates, marker loss or excessive deck motion require a go-around or safe alternative within the energy reserve. These gates and tolerances require platform-specific validation; returning to the original launch coordinates would miss the moving UGV.


%TC:ignore
\begin{table}[htbp]
\centering
\caption{Module-to-method mapping for the proposed delivery design. Sources for the named methods appear in the main text; delivery interlocks and integration rules are proposed choices.}
\label{tab:modules}
\ReplaceableGraphic{tab2_modules}{{\small\renewcommand{\arraystretch}{1.22}
\begin{tabular}{>{\raggedright\arraybackslash}p{0.24\linewidth}>{\raggedright\arraybackslash}p{0.34\linewidth}>{\raggedright\arraybackslash}p{0.32\linewidth}}
\toprule
\textbf{Module} & \textbf{Required feature} & \textbf{Example method} \\
\midrule
Hardware and sensors & Loaded flight, controlled deposition and physical confirmation. & Winch with load/contact release interlock. \\
Perception/data processing & Candidates, obstacles, target identity and evidence. & Adapted YOLO plus registered identifier check. \\
Positioning/SLAM & Vehicle state, map and deck-relative pose. & EKF; calibrated marker PnP near the deck. \\
Mission/task planning & Sequencing, action conditions and recovery choices. & Behaviour tree with bounded retry/fallback. \\
Motion planning & Collision-free, dynamically feasible timed references. & A* and constrained minimum-snap trajectory. \\
Modelling/control & Stable tracking and feasible motor commands. & Cascaded P/PID, feedforward and allocation. \\
GCS/HMI, communication and results & Mission input, vehicle updates and received delivery records. & MAVLink plus application acknowledgement. \\
\bottomrule
\end{tabular}}
}
\end{table}

\begin{figure}[htbp]
\centering
\ReplaceableGraphic{fig2_mission}{\begin{tikzpicture}[
  font=\sffamily\fontsize{8.5}{10.2}\selectfont,
  stage/.style={draw=black!65,rounded corners=2pt,fill=blue!3,
    line width=.65pt,text width=3.15cm,minimum height=1.0cm,align=center,
    inner xsep=5pt,inner ysep=5pt},
  gate/.style={diamond,aspect=1.6,draw=orange!65!black,fill=orange!8,
    line width=.7pt,text width=1.65cm,align=center,inner sep=1pt},
  note/.style={draw=black!45,rounded corners=2pt,fill=black!2,
    line width=.5pt,text width=3.15cm,minimum height=.72cm,align=center,
    inner sep=4pt},
  flow/.style={-{Latex[length=2.1mm,width=1.4mm]},line width=.8pt,
    draw=blue!50!black},
  feedback/.style={-{Latex[length=2mm,width=1.3mm]},line width=.65pt,
    draw=orange!70!black,dashed},
  edge/.style={font=\sffamily\fontsize{8}{9.5}\selectfont,
    fill=white,inner sep=2pt,align=center}
]
  \node[stage] (ready) at (0,0)
    {\textbf{Readiness checks}\\Payload, localisation,\\route and reserve};
  \node[stage] (launch) at (4.8,0)
    {\textbf{Launch from the UGV}\\Match deck velocity;\\climb clear};
  \node[stage] (search) at (9.6,0)
    {\textbf{Transit and search}\\Reach building;\\scan target fa\c{c}ade};

  \node[gate] (check) at (9.6,-2.1)
    {Correct target\\and clear\\area?};
  \node[stage] (deliver) at (4.8,-2.1)
    {\textbf{Approach and delivery}\\Stable placement;\\controlled release};
  \node[stage] (verify) at (0,-2.1)
    {\textbf{Verify delivery}\\Camera; payload sensor;\\record success/failure};

  \node[stage] (report) at (0,-4.4)
    {\textbf{Send delivery record}\\Outcome, time, target ID;\\track receipt status};
  \node[stage] (meet) at (4.8,-4.4)
    {\textbf{Update rendezvous}\\UGV pose and velocity;\\check timestamps};
  \node[stage] (recover) at (9.6,-4.4)
    {\textbf{Track and recover}\\Align; match velocity;\\descend to contact};

  \draw[flow] (ready) -- (launch);
  \draw[flow] (launch) -- (search);
  \draw[flow] (search) -- (check);
  \draw[flow] (check) -- node[edge,above]{Yes} (deliver);
  \draw[flow] (deliver) -- (verify);
  \draw[flow] (verify) -- (report);
  \draw[flow] (report) -- node[edge,above]{Return} (meet);
  \draw[flow] (meet) -- (recover);

  \coordinate (retryfork) at (11.85,-2.1);
  \draw[feedback] (check.east) -- node[edge,above]{No} (retryfork);
  \draw[feedback] (retryfork) -- (11.85,.95)
    -- node[edge,above]{Retry within budget} (9.6,.95) -- (search.north);

  \node[note] (receipt) at (0,-5.85)
    {No ACK: retain report;\\retry within budget.};
  \draw[feedback,draw=black!55] (report.south) --
    node[edge,right]{In parallel} (receipt.north);

  \node[note] (around) at (9.6,-5.85)
    {Unsafe: go around\\if energy reserve permits.};
  \draw[feedback] (recover.south) -- (around.north);
  \draw[feedback] (around.west) -- (7.2,-5.85) -- (7.2,-5.45)
    -- node[edge,above]{Retry} (4.8,-5.45) -- (meet.south);

  \node[note,text width=12.6cm,minimum height=.7cm,fill=orange!5,
        draw=orange!65!black] (fallback) at (4.8,-7.2)
    {\textbf{Safety fallback:} cancel release; return safely or use a designated alternative landing site.};
  \draw[feedback] (retryfork) -- (11.85,-7.2) -- (fallback.east);
  \node[edge,text width=1.2cm] at (12.6,-3.1) {Budget\\exhausted};
  \draw[feedback] (around.south) --
    node[edge,right]{Otherwise} (9.6,-6.82);

  \node[align=left,font=\sffamily\fontsize{8}{9.5}\selectfont,
        text width=3.8cm] at (4.8,-6.15)
    {Report receipt and physical\\delivery are separate states.};
\end{tikzpicture}
}
\caption{Proposed mission flow with verification and fallback conditions. The task logic uses behaviour-tree principles \citep{Colledanchise2018}; this diagram is a mission flowchart. Physical delivery and reception of its report are separate states.}
\label{fig:mission}
\end{figure}

\begin{figure}[htbp]
\centering
\ReplaceableGraphic{fig3_recovery}{\begin{tikzpicture}[
  font=\sffamily\fontsize{8.5}{10.2}\selectfont,
  arr/.style={-{Latex[length=2.1mm,width=1.4mm]},line width=.8pt},
  dim/.style={{Latex[length=1.6mm,width=1.1mm]}-{Latex[length=1.6mm,width=1.1mm]},
    line width=.65pt,draw=black!65},
  lab/.style={align=center,fill=white,inner sep=2pt},
  faint/.style={draw=black!45,dashed,line width=.65pt}
]
  \node[font=\sffamily\bfseries\fontsize{9}{11}\selectfont]
    at (4.65,6.25) {Side view: level deck with straight motion};
  \node[align=center] at (4.9,5.65)
    {Track the moving deck while descending\\
     $\mathbf v_{U,xy}\approx\mathbf v_{D,xy}$};

  % A single z-up world frame is used by both vehicle states.
  \draw[black!35,line width=.65pt] (.25,.6) -- (9.1,.6);
  \draw[arr] (.6,.6) -- (1.8,.6) node[right] {$x$};
  \draw[arr] (.6,.6) -- (.6,1.8) node[above] {$z$};
  \node[anchor=west,align=left,text width=2.5cm] at (.3,-.12)
    {World frame: $z$ up};

  % Moving UGV and its recovery deck.
  \draw[draw=black!65,fill=blue!5,line width=.8pt]
    (2.5,1.15) rectangle (7.7,1.8);
  \draw[draw=black!65,fill=blue!5,line width=.8pt]
    (7.7,1.15) -- (8.5,1.15) -- (8.5,1.5) -- (8.1,1.8) -- (7.7,1.8) -- cycle;
  \draw[black!65] (7.85,1.7) -- (8.05,1.7) -- (8.3,1.5) -- (7.85,1.5) -- cycle;
  \foreach \xx in {3.1,7.75}{
    \draw[draw=black!65,fill=white,line width=.8pt] (\xx,.91) circle (.27);
    \fill[black!55] (\xx,.91) circle (.055);
  }
  \node at (5.1,1.42) {Moving UGV};
  \draw[draw=blue!55!black,fill=blue!15,line width=1pt]
    (4.35,1.79) rectangle (5.45,1.85);
  \fill[blue!65!black] (4.9,1.85) circle (.055);
  \node[anchor=west,lab] at (5.55,2.05) {$\mathbf p_D(t)$: deck target};
  \draw[arr,draw=blue!65!black] (7.75,2.15) -- (9.05,2.15)
    node[midway,above] {$\mathbf v_{D,xy}$};

  % The UAV glyph is centred at its centre of mass; feet lie 0.75 units below it.
  \def\recoveryuav{%
    \draw[line width=.8pt] (-.82,0) -- (.82,0);
    \draw[fill=white,line width=.8pt] (-.38,-.13) rectangle (.38,.13);
    \draw[line width=.8pt] (-.8,.04) ellipse (.34 and .065);
    \draw[line width=.8pt] (.8,.04) ellipse (.34 and .065);
    \draw[line width=.8pt] (-.27,-.13) -- (-.45,-.75) -- (-.7,-.75);
    \draw[line width=.8pt] (.27,-.13) -- (.45,-.75) -- (.7,-.75);
  }
  \begin{scope}[shift={(4.9,4.8)},draw=blue!65!black]
    \recoveryuav
    \fill[blue!65!black] (0,0) circle (.055);
  \end{scope}
  \node[lab] at (4.9,5.15) {UAV COM: $\mathbf p_U(t)$};
  \draw[arr,draw=blue!65!black] (6.2,4.8) -- (8.6,4.8)
    node[midway,above] {$\mathbf v_{U,xy}$};

  \begin{scope}[shift={(4.9,2.6)},draw=black!45,dashed]
    \recoveryuav
    \fill[black!50] (0,0) circle (.055);
  \end{scope}
  \draw[faint] (4.9,4.63) -- (4.9,2.75);
  \draw[arr,draw=orange!75!black] (5.7,3.95) -- (5.7,3.0);
  \node[anchor=west,align=left,text width=2.8cm] at (6.0,3.55)
    {Controlled descent;\\vertical relative speed\\approaches zero\\at touchdown};

  % Heights refer to centre of mass, not to the landing-gear contact point.
  \draw[faint] (2.15,4.8) -- (3.35,4.8);
  \draw[faint] (2.15,1.85) -- (4.25,1.85);
  \draw[dim] (2.2,1.85) -- node[left,lab] {$h(t)$} (2.2,4.8);
  \draw[faint] (3.25,2.6) -- (3.7,2.6);
  \draw[dim] (3.3,1.85) -- node[left,lab] {$h_{\rm td}>0$} (3.3,2.6);
  \node[align=center,text width=5.8cm] at (6.1,-.12)
    {$h(t)\downarrow h_{\rm td}$, the COM clearance at contact.};

  % Compact top view makes the horizontal alignment condition explicit.
  \draw[draw=black!35,rounded corners=3pt,line width=.6pt,fill=black!1]
    (9.6,.55) rectangle (14.95,6.55);
  \node[font=\sffamily\bfseries\fontsize{9}{11}\selectfont]
    at (12.275,6.25) {Top view: horizontal alignment};
  \draw[draw=black!55,fill=blue!5,line width=.7pt]
    (10.35,3.7) rectangle (14.2,5.35);
  \draw[faint] (12.15,4.5) circle (.49);
  \draw[draw=blue!65!black,line width=.9pt]
    (11.98,4.5) -- (12.32,4.5) (12.15,4.33) -- (12.15,4.67);
  \fill[blue!65!black] (12.15,4.5) circle (.04);
  \fill[orange!80!black] (12.55,4.76) circle (.06);
  \draw[arr,draw=orange!80!black] (12.17,4.52) -- (12.5,4.73);
  \node[anchor=east,lab] at (12.02,5.03) {$\mathbf e_{xy}$};
  \draw[arr,draw=blue!65!black] (12.25,4.35) -- (13.85,4.35);
  \draw[arr,draw=orange!80!black] (12.65,4.98) -- (14.25,4.98);
  \node[font=\sffamily\fontsize{8}{9.5}\selectfont,align=center,
        text width=4.7cm] at (12.28,3.35)
    {Cross: deck target; dot: UAV.\\Arrows: horizontal velocity.};

  \node at (12.275,2.8) {\textbf{Permit descent only if}};
  \node at (12.275,2.35) {$\|\mathbf e_{xy}\|\leq\varepsilon_p$};
  \node at (12.275,1.9)
    {$\|\mathbf v_{U,xy}-\mathbf v_{D,xy}\|\leq\varepsilon_v$};
  \node at (12.275,1.4) {Reliable pose; clear deck.};
  \node[align=center,font=\sffamily\fontsize{8}{9.5}\selectfont,
        text width=4.6cm] at (12.275,.88)
    {$U$: UAV COM; $D$: deck target.};
\end{tikzpicture}
}
\caption{Proposed moving-deck recovery in a common, locally level frame. Mobile-target tracking and camera-pose estimation draw on \citet{Borowczyk2017} and \citet{Lepetit2009}. Horizontal matching and a tapered touchdown speed are design choices for this scenario; the touchdown height of the centre of mass remains positive.}
\label{fig:recovery}
\end{figure}
%TC:endignore
